% TOP_PROBLEM: {"id": 1123, "title": "The Illumination Conjecture", "category_id": 6, "category": {"id": 6, "name": "geometry", "display_name": "Geometry", "description": "Euclidean and non-Euclidean geometry, geometric structures.", "slug": "geometry", "order_index": 6, "created_at": "2026-07-31T15:26:25.671Z"}, "status": "open"}
% FIRST_POSTED: 2026-09-27
% !TeX program = lualatex
% ATTEMPT_STATUS: unresolved
\documentclass[11pt]{article}
\usepackage[margin=25mm]{geometry}
\usepackage{fontspec}
\setmainfont{Latin Modern Roman}
\usepackage{amsmath,amssymb,mathtools}
\usepackage{unicode-math}
\setmathfont{Latin Modern Math}
\usepackage{xurl,hyperref}
\hypersetup{colorlinks=true,urlcolor=blue}
\setcounter{secnumdepth}{0}
\setlength{\parindent}{0pt}
\setlength{\parskip}{5pt}
\setlength{\emergencystretch}{3em}
\begin{document}
\section*{\texorpdfstring{The Illumination Conjecture}{The Illumination Conjecture}}\label{1123}
\subsection{Definitions and mathematical statement}
A convex body $K\subset{\mathbb{R}}^n$ is compact and convex with nonempty interior. A direction $v\in S^{n-1}$ illuminates $x\in\partial K$ if $x+tv\in\operatorname{int}K$ for some $t>0$. Its illumination number is
\[
 I(K)=\min\{|D|:D\subset S^{n-1},\
 \forall x\in\partial K\ \exists v\in D\ \exists t>0:\
 x+tv\in\operatorname{int}K\}.
\]
The conjecture is $I(K)\le2^n$ for every dimension and every convex body. The equivalent external-point-source formulation lets a ray from a source continue through the boundary point into the interior. One direction may illuminate many boundary points; the finite collection must cover all of the boundary, not merely almost every point.
\par\smallskip\textit{Formulation and concept references:} \href{https://arxiv.org/pdf/2507.08712}{[S1]}.

\subsection{Short English statement}
Can the boundary of every n-dimensional convex body be illuminated using at most two to the n light directions?
\subsection{Research attempt}
\paragraph{Scope and source check (27 September 2026).}
The target concerns every convex body and every boundary point. The cap-body theorem in [S1] is a result for a particular class in dimension three: the checked version proves an upper bound of six for that class. It does not establish the stated all-dimensional assertion. This attempt proves the conjectured bound for smooth bodies, simplices, parallelotopes, and products of bodies already covered by a bound. It also identifies exactly why approximation by smooth bodies cannot complete the argument. All proofs below concern the directional definition in the statement; a tangent direction does not count as illumination.

\paragraph{Lemma 473.1: an exact normal-cone criterion.}
For $x\in\partial K$ let
\[
 N_K(x)=\{u\in S^{n-1}:\langle u,z-x\rangle\leq0\text{ for all }z\in K\}.
\]
This is the compact nonempty set of unit outward supporting normals at $x$. A direction $v$ illuminates $x$ if and only if
\begin{equation}\label{a473:normal}
       \langle u,v\rangle<0\quad\text{for every }u\in N_K(x).
\end{equation}
Necessity follows because an interior point lies strictly inside every supporting half-space. Here is a proof of sufficiency that also treats nonsmooth points. If no sufficiently small positive step enters the interior, take $t_j\downarrow0$ with $x+t_jv\notin\operatorname{int}K$. Separation gives a unit vector $u_j$ with
\[
 \langle u_j,x+t_jv\rangle\geq h_K(u_j),\qquad
 h_K(u)=\max_{z\in K}\langle u,z\rangle.
\]
Since $\langle u_j,x\rangle\leq h_K(u_j)$, we have $\langle u_j,v\rangle\geq0$. Passing to a convergent subsequence and using continuity of the support function gives a limit $u\in N_K(x)$ with $\langle u,v\rangle\geq0$, contradicting \eqref{a473:normal}. In particular the strict criterion makes all sufficiently small steps work, not just one isolated step.

For a fixed point, compactness of $N_K(x)$ turns the strict inequalities into a margin $\langle u,v\rangle\leq-\delta$ for some $\delta>0$. That margin depends on the point and body. Losing it in a limit is a substantive obstruction, as shown below.

\paragraph{Finite illumination and the universal lower bound.}
Translate an interior point to the origin. For every boundary point $x$, the direction $-x/\|x\|$ illuminates $x$. If $x+tv$ is interior, then $y+tv$ stays interior for all $y$ sufficiently close to $x$. Thus each direction illuminates a relatively open subset of the boundary. Compactness produces a finite subcover, proving that $I(K)$ is finite, with no uniform count.

On the other hand,
\begin{equation}\label{a473:lower}
                         I(K)\geq n+1.
\end{equation}
Indeed, if at most $n$ vectors form $D$, then zero is not in the interior of their convex hull in $\mathbb R^n$. A separating hyperplane gives a nonzero $u$ with $\langle u,v\rangle\geq0$ for all $v\in D$. Choose a support point of $K$ with outward normal $u/\|u\|$. Criterion \eqref{a473:normal} excludes every direction in $D$ at that point. This uses a support point for the chosen normal, not an assumption that the boundary has a unique normal.

\paragraph{Proposition 473.2: smooth bodies need exactly $n+1$ directions.}
Suppose that every boundary point has a unique supporting unit normal. Choose the $n+1$ unit directions of a regular simplex centered at zero. They span the space and have positive coefficients, for instance equal ones, summing to zero:
\[
                    v_0+\cdots+v_n=0.
\]
For any nonzero $u$, at least one $\langle u,v_i\rangle$ is negative. Otherwise their sum is zero with every summand nonnegative, so all are zero, contrary to spanning. At a smooth boundary point there is just one normal to test, and that negative direction illuminates it. Hence $I(K)\leq n+1$, and \eqref{a473:lower} proves equality.

It is not enough to apply this reasoning separately to each normal at a corner: the direction must be negative on all of that point's normals simultaneously. Changing the direction when changing the normal reverses the required order of the quantifiers. This is the precise place at which the smooth proof stops.

\paragraph{Affine invariance and a polytope certificate.}
If $T(x)=Ax+b$ with $A$ invertible, an illuminating direction $v$ becomes $Av/\|Av\|$ for $T(K)$. The parameter of the ray is merely rescaled by the positive number $\|Av\|$. Applying the inverse affine map proves
\[
                              I(T(K))=I(K).
\]
For a full-dimensional polytope, it suffices to illuminate all vertices. Write the polytope using its facet inequalities. At a point $x$ of a face, take any vertex $w$ of that face. Every inequality active at $x$ is also active at $w$. A vector strictly decreasing all active facet functionals at $w$ therefore decreases all those active at $x$, and a sufficiently small step keeps the remaining, strict inequalities satisfied. Thus any collection covering the vertices covers the entire boundary.

For a simplex, choose at each vertex the vector pointing to an interior point. These $n+1$ directions cover its vertices, so its illumination number is exactly $n+1$. No regularity of the simplex is needed, either by this argument or by affine invariance.

For the cube $Q=[-1,1]^n$, a vertex $\varepsilon\in\{-1,1\}^n$ is illuminated precisely when
\begin{equation}\label{a473:cube}
                         \varepsilon_i v_i<0\quad(1\leq i\leq n).
\end{equation}
One vector cannot satisfy this for two different vertices: at a differing coordinate it would need both signs. Hence $I(Q)\geq2^n$. The $2^n$ opposite sign vectors, normalized to unit length, do illuminate all vertices. Therefore
\[
                       I(Q)=2^n,
\]
and affine invariance proves the same equality for every parallelotope. This supplies a sharp example for the conjectured upper bound, rather than a reason to expect a smaller universal bound.

\paragraph{Proposition 473.3: the product construction.}
For convex bodies $K\subset\mathbb R^a$ and $L\subset\mathbb R^b$,
\begin{equation}\label{a473:product}
                           I(K\times L)\leq I(K)I(L).
\end{equation}
Take illuminating sets $D_K,D_L$, and use the normalized vectors $(v,w)$ for all $(v,w)\in D_K\times D_L$. For a boundary point $(x,y)$, if $x$ is a boundary point choose $v$ illuminating $x$; if $x$ is interior choose any $v$ in the set. Make the corresponding choice for $y$. In each coordinate an entire sufficiently short positive segment stays interior after leaving the boundary, or stays interior from the outset. Choosing a common sufficiently small parameter puts both coordinates in the interior. Normalization does not change this property. Consequently a product of bodies obeying their respective $2^{\dim}$ bounds also obeys the desired bound. This proves the bound for all products of the special cases above.

\paragraph{The combinatorial form of the missing polytope argument.}
For a subset $V$ of the vertices of a polytope, collect the unit normals of all facets active at some member of $V$ into a finite set $A_V$. There is one direction illuminating every vertex in $V$ if and only if
\begin{equation}\label{a473:cluster}
                          0\notin\operatorname{conv}(A_V).
\end{equation}
One implication follows by taking the inner product of a convex combination with a direction strictly negative on every member. Conversely, strict separation of the compact convex hull from zero supplies such a direction. Homogeneity allows it to be normalized. With rational facet data the test can instead use unnormalized rational normals and the linear inequalities $\langle a,v\rangle\leq-1$, which are feasible exactly when a strict solution exists after rescaling.

For a specified finite polytope this gives a verifiable certificate: partition its vertices into compatible subsets satisfying \eqref{a473:cluster}, and record one feasible vector for each subset. The conjecture would follow for that polytope if the partition had at most $2^n$ parts. Merely having an algorithm to test compatibility does not prove that such a partition always exists. The number of vertices is unrestricted, and the elementary argument only gives one direction per vertex.

There is also no fixed universal list of $2^n$ directions that can simply be chosen independently of the body. Already in the plane, choose an inward direction missing a given finite list and a sufficiently narrow open cone around it that misses the whole list. A triangle with that inward cone at one vertex defeats the list. This does not defeat the conjecture, because the conjecture permits choosing directions after seeing the body.

A less extreme check is the diamond $|x_1|+|x_2|\leq1$. The four diagonal directions $(\pm1,\pm1)$, even though they solve the axis-aligned square, do not illuminate $(1,0)$ on this diamond. At that vertex the active normals are $(1,1)$ and $(1,-1)$. Each leftward diagonal has zero inner product with one of them, and each rightward diagonal has a positive one. Replacing strict negativity by nonpositivity would incorrectly mark this example as covered.

\paragraph{A decisive countercheck to smoothing and compactness.}
For $p>1$ put
\[
                 K_p=\{x:\ |x_1|^p+\cdots+|x_n|^p\leq1\}.
\]
$K_p$ has a $C^1$ boundary and a unique supporting normal at each boundary point, so the proved smooth case gives $I(K_p)=n+1$. Moreover
\[
                       n^{-1/p}Q\subset K_p\subset Q.
\]
These inclusions imply Hausdorff convergence $K_p\longrightarrow Q$ as $p\to\infty$. But $I(Q)=2^n$. For $n\geq2$ this is strictly larger than $n+1$.

Thus an implication of the form ``smooth approximants admit $m$ directions, therefore the limit admits $m$ directions'' is false, even in this explicit symmetric family. Extracting convergent subsequences of the illuminating directions only gives limiting vectors; the step sizes or the strict normal margins can vanish. At a limiting cube vertex one must satisfy all coordinate inequalities \eqref{a473:cube} at once, whereas each approximating smooth point has only one supporting normal. A successful approximation method needs an additional uniform strictness or a controlled enlargement of the direction set. Neither follows from compactness alone.

\paragraph{Finite certificates and their limits.}
The accompanying exact diagnostic uses integer inner products to verify the cube sign obstruction in dimensions one through eight, the diamond failure above, and the simplex-direction property on a bounded set of nonzero integer normals. The simplex vectors are represented without normalization by $e_1,\ldots,e_n,-\sum e_i$: their positive dependence and spanning give the proof for every normal, while the enumeration only checks the implementation of that criterion. The product proof and the normal-cone separation argument are mathematical deductions, not consequences of a numerical plot. No floating-point tolerance is used to turn tangency into illumination.

These calculations validate concrete certificates and failed shortcuts. They do not enumerate all convex bodies, all facet configurations, or all dimensions. In particular, the finite compatibility criterion is not accompanied here by a theorem limiting the required number of compatible groups.

\paragraph{What would close the gap.}
The main unresolved step is a uniform covering theorem for the normal cones of an arbitrary convex body. For polytopes one precise sufficient statement is that their vertices admit a partition into at most $2^n$ sets satisfying \eqref{a473:cluster}. A route through approximation must additionally preserve strict inwardness at every limiting boundary point; the $K_p$ example shows why this requirement cannot be suppressed. The known cap-body theorem in [S1] addresses a special geometric configuration, and the arguments here establish several other special configurations. They do not provide this general covering theorem. The all-body conjecture therefore remains unresolved in this attempt.

\textbf{UNRESOLVED.} The partial results above do not establish the full selected statement.

\subsection{Sources}
\begin{itemize}
\item[S1]Arman, A., Kaire, J.S., Prymak, A., 'Illumination number of 3-dimensional cap bodies', Discrete Math. 349(7) (2026); arXiv:2507.08712.
\url{https://arxiv.org/pdf/2507.08712}.
\textit{Formulation source.}
\end{itemize}
\end{document}
